\chapter{Клеточное типизированное вложение
\texorpdfstring{$SU(2)$}{SU(2)}-посредника}
\label{ch:v10-discrete-resolution-su2-mediator-cell-embedding}

Пусть $B=(-1,1)^\top$ --- граница одного ориентированного ребра. Поле
посредника несёт adjoint-представление, поэтому его клеточная граница и
лапласиан равны
\begin{equation}
 B_{\rm ad}=B\otimes I_3,\qquad
 \Delta_{\rm ad}=B_{\rm ad}B_{\rm ad}^\top
 =2P_\perp\otimes I_3.
\end{equation}
Они имеют ранги $3$ и ядро размерности $3$: по одной постоянной gauge-моде
для каждой компоненты алгебры. Замена ориентации ребра не меняет лапласиан.

Теперь типизируем источник. Для двух дублетов введём полные и относительные
генераторы
\begin{equation}
 J_a=T_a\otimes I+I\otimes T_a,\qquad
 R_a=T_a\otimes I-I\otimes T_a.
\end{equation}
Полные генераторы аннулируют синглет $|S\rangle$, тогда как три вектора
$R_a|S\rangle$ ортонормальны, натягивают весь триплет и преобразуются по
adjoint-представлению. Матрица
\begin{equation}
 V=(R_1|S\rangle,R_2|S\rangle,R_3|S\rangle)
\end{equation}
удовлетворяет $V^\dagger V=I_3$ и $VV^\dagger=P_T$.

Пространственное вложение $E=B_{\rm ad}/\sqrt2$ также изометрично и лежит
в сохранённом секторе. Поэтому
\begin{equation}
 C=EV^\dagger:\mathbb C^4\longrightarrow\mathbb C^6
\end{equation}
является частичной изометрией с $C^\dagger C=P_T$ и
$CC^\dagger=P_\perp\otimes I_3$. Проецированный Green-оператор даёт
\begin{equation}
 E^\dagger\left(\frac12P_\perp\otimes I_3\right)E
 =\frac12I_3,
\end{equation}
то есть все три adjoint-компоненты имеют одинаковую восприимчивость и
сохраняют условный коэффициент $\kappa=3/16$.

\begin{theorem}[Клеточный adjoint-ток типизирован]
Относительные токи композитной пары и разностные клеточные токи связаны
точной $SU(2)$-ковариантной частичной изометрией ранга три. Нулевая gauge-
мода не возбуждается, а статический отклик одинаков по всем компонентам.
\end{theorem}

\begin{nogocriterion}[Типизация не выбирает flavor-пару]
Построение действует одинаково на все шестнадцать копий дублета. Оно не
выбирает одну из $256$ flavor-пар и не связывает рациональный gap $3/16$ с
трансмутированной массой полюса.
\end{nogocriterion}

\begin{protocol}\begin{enumerate}
\item условная архитектура: \textbf{$14/14$};
\item типизированное вложение: \textbf{$10/10$};
\item унаследованные ингредиенты: \textbf{$2/2$};
\item flavor-селектор и pole-map: \textbf{$0/2$};
\item следующий гейт: \textbf{аудит селектора flavor-пары}.
\end{enumerate}\end{protocol}

Машинный сертификат сохранён в
\path{s2t_v10_cell_birth_four_volume_nonequilibrium_bath_discrete_resolution_transmuted_mode_asymptotically_free_su2_gauge_singlet_su2_mediator_cell_complex_typed_embedding_gate_results.json}.